%% ma: L10-B13-0012
%% lop: 10
%% chuong: 5
%% bai: 13
%% dang: 10.13.2
%% loai: TN4
%% muc: TH
%% dap_an: "C"
%% nguon: "chuyen-de-cac-so-dac-trung-cua-mau-so-lieu-khong-ghep-nhom-toan-10-chuong-trinh-moi.pdf (GV. Phan Nhật Linh) (P02), Bài 02 Dạng 1 Phần I Câu 2"
%% kien_thuc: "Lớp 10 Bài 13 (số trung bình)"
%% trang_thai: chinh-thuc
%% ghi_chu: "Phương án C gốc bị trống, đã điền 3,7 theo lời giải. Tự duyệt theo yêu cầu giáo viên 22/10/2026."
\begin{ex}
Số nhân khẩu trong các hộ gia đình ở một xóm được thống kê ở bảng sau:
\begin{center}
\begin{tabular}{|c|c|c|c|c|c|c|}\hline
Số nhân khẩu & $1$ & $2$ & $3$ & $4$ & $5$ & $6$\\\hline
Số hộ gia đình & $1$ & $4$ & $7$ & $11$ & $5$ & $2$\\\hline
\end{tabular}
\end{center}
Số trung bình của mẫu số liệu trên là
\choice{$3{,}5$}{$2$}{\True $3{,}7$}{$5$}
\loigiai{$n=30$; $\overline x=\dfrac{1\cdot1+2\cdot4+3\cdot7+4\cdot11+5\cdot5+6\cdot2}{30}=\dfrac{111}{30}=3{,}7$. Chọn C.}
\end{ex}
